\subsection{Partial derivatives}

1.6.1. Let \(f\) be a function defined in a neighborhood \(U\) of the point \(x_0\) of \(E\) and with values in \(F\) . Let \(X\) be a vector subspace of \(E\) and \(V\) the set of points \(x\) of \(X\) such that \(x_0 + x \in U\) ; set \(g(x) = f(x_0 + x)\) for \(x \in V\) . We say that \(f\) admits a partial derivative with respect to \(X\) at \(x_0\) if \(g\) admits a derivative at 0; this derivative is denoted \(D_X f(x_0)\) ; it is a continuous linear map from \(X\) to \(F\) . If \(f\) is differentiable at \(x_0\) it admits a partial derivative with respect to \(X\) at \(x_0\) and this partial derivative is the restriction of \(Df(x_0)\) to \(X\) .

1.6.2. Suppose that \(E\) is the product of a finite family of normed vector spaces \(E_{i}\) ( \(1 \leqslant i \leqslant n\) ) canonically identified with subspaces of \(E\); let \(x_{0} = (x_{0}^{1}, \ldots, x_{0}^{n})\) be in \(E\) and let \(U\) be a neighborhood of \(x_{0}\) in \(E\); finally let \(f\) be a mapping from \(U\) into \(F\). We denote by \(D_{i}f(x_{0})\) the derivative at the point \(x_{0}^{i}\) , if it exists, of the mapping \(z_{i} \mapsto f(x_{0}^{1}, \ldots, z_{i}, \ldots, x_{0}^{n})\) defined in a neighborhood of \(x_{0}^{i}\) in \(E_{i}\) and taking values in \(F\). This is an element of \(\mathcal{L}(E_{i}; F)\) which is called the i-th partial derivative of \(f\) at \(x_{0}\) . If \(f\) is differentiable at \(x_{0}\) , the \(n\) partial derivatives exist, and determine \(Df(x_{0})\) by the formula:

\[
\mathrm{D} f (x _ {0}). h = \sum_ {i = 1} ^ {n} \mathrm{D} _ {i} f (x _ {0}). h _ {i} \quad \text { for } h = (h _ {1}, \dots , h _ {n}) \text { in } E.
\]

1.6.3. More particularly, let \(E = K^{n}\) . If the partial derivatives of f at \(x_{0}\) exist, we denote \(\partial_{i}f(x_{0})\) the element \(D_{i}f(x_{0}).1\) of F. The following notation is often used; suppose that a notation has been chosen for the coordinate functions on \(K^{n}\) , for example \(u_{i}\) denotes the i-th projection of \(K^{n}\) onto K. One then writes:

\[
\frac {\partial f}{\partial u _ {i}} \left(x _ {0}\right) \quad \text { or } \quad \left. \frac {\partial f}{\partial u _ {i}} \right| _ {x = x _ {0}}
\]

instead of \(\partial_{i}f(x_{0})\) .

1.6.4. Let \(E = K^{n}\) and \(F = K^{m}\) ; suppose that the function \(f = (f_{1}, \ldots, f_{m})\) with values in F is differentiable at the point \(x_{0}\) of E. The partial derivatives \(a_{ji} = \partial_{i} f_{j}(x_{0})\) then exist (they are elements of K). The matrix with m rows and n columns formed by the \(a_{ji}\) (element in the row of index j and the column of index i) is called the Jacobian matrix of f at \(x_{0}\) ; this is the matrix of the linear map \(Df(x_{0})\) of \(K^{n}\) to \(K^{m}\) with respect to the canonical bases of these spaces.

